\documentclass[11pt,twocolumn]{article}
\usepackage[utf8]{inputenc}
\usepackage{amsmath,amssymb,amsthm}
\usepackage{graphicx}
\usepackage{algorithm}
\usepackage{algpseudocode}
\usepackage{booktabs}
\usepackage{xcolor}
\usepackage{hyperref}
\usepackage[margin=1in]{geometry}
\usepackage{listings}
\usepackage{verbatim}

\title{\textbf{Automatic Rule Concept Surrogate for Deep Reinforcement Learning}}
\author{Anonymous Authors}
\date{}

\begin{document}

\maketitle

\begin{abstract}
Interpreting deep reinforcement learning (RL) policies remains a fundamental challenge in understanding agent behavior. We introduce \textbf{ARC} (Automatic Rule Concept surrogate), a framework that automatically extracts human-interpretable boolean rules from trained deep RL policies while maintaining high task performance (90--100\% success rate). Unlike prior methods that require manual concept specification or sacrifice accuracy for interpretability, ARC combines sparse autoencoders with product t-norm fuzzy logic to discover both the concepts and the logical rules governing agent decisions. Our two-stage training approach prevents gradient conflicts and produces concise, verifiable rules such as ``Forward $\leftarrow$ ($\lnot f_{25} \land f_{26} \land \lnot f_{81} \land \lnot f_{91} \land \lnot f_{280}$)'' that match or exceed the original policy's performance on MiniGrid environments.
\end{abstract}

\section{Introduction}

\subsection{Motivation}

Deep reinforcement learning has achieved remarkable success across diverse domains \cite{silver2016alphago,mnih2015human}. However, the black-box nature of neural policies hinders their deployment in safety-critical applications and limits our understanding of learned behaviors. When an agent makes a decision, practitioners need answers to fundamental questions: \emph{Why did the agent choose this action? What perceptual features drove this decision? Can we trust this policy in novel situations?}

Interpretability in deep RL requires more than post-hoc explanations—it demands extractable, verifiable rules that faithfully represent the policy's decision logic. Such rules enable formal verification, facilitate debugging, and build trust through transparency.

\subsection{Limitations of Existing Approaches}

Current interpretability methods face significant limitations:

\paragraph{Concept Bottleneck Models} \cite{koh2020concept} achieve interpretability by routing decisions through pre-defined human concepts. However, they require \emph{manual specification} of concepts before training, making them:
\begin{itemize}
\item \textbf{Labor-intensive:} Domain experts must enumerate all relevant concepts a priori
\item \textbf{Incomplete:} Hand-crafted concepts may miss emergent features the policy discovers
\item \textbf{Rule-agnostic:} Even with concepts, extracting logical rules (especially complex conjunctions) remains difficult
\end{itemize}

\paragraph{Decision Trees and Rule Extraction} Methods like CART \cite{breiman1984classification} and soft decision trees \cite{frosst2017distilling} offer rule-based interpretability but suffer from:
\begin{itemize}
\item \textbf{Binary splits:} Each node splits on a single feature, creating deep trees for multi-variate rules
\item \textbf{Accuracy degradation:} Shallow trees lose expressiveness; deep trees become uninterpretable hierarchies
\item \textbf{Limited expressiveness:} Struggle to represent conjunctive rules like ``Pickup $\leftarrow$ ($f_1 \land f_2 \land \lnot f_3 \land f_4$)'' compactly
\end{itemize}

\subsection{Our Contribution: ARC Framework}

We propose ARC, which automatically:
\begin{enumerate}
\item \textbf{Discovers concepts} from raw observations via overcomplete sparse autoencoders, eliminating manual specification
\item \textbf{Learns boolean rules} in Disjunctive Normal Form (DNF) using product t-norm fuzzy logic, enabling compact multi-literal conjunctions
\item \textbf{Maintains high accuracy} (90--100\% success rate on MiniGrid-DoorKey) through a two-stage training protocol that prevents gradient interference
\end{enumerate}

The key insight is a \emph{separation of representation learning from rule learning}: freeze the sparse autoencoder after convergence, then train the logic layer on stable features. This architectural constraint eliminates gradient conflicts that plagued prior joint training approaches.

\section{Method}

\subsection{Problem Formulation}

Given a trained policy $\pi_\text{PPO}$ with frozen feature extractor $\phi: \mathcal{O} \to \mathbb{R}^d$ (e.g., CNN), our goal is to learn an interpretable surrogate $\pi_\text{ARC}$ such that:
\begin{align}
\pi_\text{ARC}(o) &\approx \pi_\text{PPO}(o) \quad \forall o \in \mathcal{O} \\
\pi_\text{ARC} &= \text{argmax}_a \, \bigvee_{i=1}^{n_a} C_{a,i}(c(o))
\end{align}
where $c: \mathbb{R}^d \to \{0,1\}^D$ maps features to binary concepts and $C_{a,i}$ are conjunctive clauses (AND of literals). Each action $a$ is selected by a DNF rule: OR of conjunctions.

\subsection{Architecture Overview}

ARC consists of three modules operating sequentially (Figure~\ref{fig:architecture}):

\begin{figure*}
\centering
\small
\begin{verbatim}
Raw observation --> CNN (frozen) --> Stage 1: SVD+ICA --> Normalize
                                         |                    |
                                    V_k (signal subspace)     |
                                         |                    v
                          Stage 2A: SAE Pre-training    Features (N, 128)
                          (reconstruction only)              |
                                         |                    v
                          SAE Encoder: Linear(128 -> 300) + TopK(k=50)
                                         |                    |
                          SAE Decoder: Linear(300 -> 128)    |
                          [frozen after convergence]          |
                                                              v
                          Sparse z (N, 300), k=50 active per sample
                                         |
                                         v
                          Fixed Normalization: (z - z_mean) / z_std
                                         |
                                         v
                          Stage 2B: Logic Training (SAE frozen)
                                         |
                          Sigmoid Bottleneck: sigmoid(alpha(z - beta))
                          [pushes to {0,1} via bimodality loss]
                                         |
                                         v
                          Product T-Norm Logic Layer (DNF)
                          Each action: OR(clause_1, ..., clause_n)
                          Each clause: AND over soft literals
                          literal = p*f + n*(1-f) + (1-p-n)*1
                          [p, n from softmax]
                                         |
                                         v
                          Action logits (N, 7)
\end{verbatim}
\caption{ARC architecture. Two-stage training prevents gradient wars: SAE learns stable concepts (Stage 2A), then logic layer learns rules on frozen features (Stage 2B).}
\label{fig:architecture}
\end{figure*}

\paragraph{Stage 1: Feature Space Analysis (Preprocessing)}
Given $N$ rollout samples $\{x_i, a_i\}_{i=1}^N$ where $x_i = \phi(o_i) \in \mathbb{R}^d$:
\begin{itemize}
\item \textbf{SVD}: Decompose $X \in \mathbb{R}^{N \times d}$ to identify signal subspace $V_k$ capturing 95\% variance
\item \textbf{ICA}: Extract $k$ independent components as initialization anchors for SAE decoder
\item \textbf{Normalization}: Compute $\mu_x, \sigma_x$ for feature standardization
\end{itemize}

This stage provides: (1) stable initialization for SAE via ICA directions, (2) signal subspace $V_k$ for regularization, (3) normalization statistics.

\paragraph{Stage 2A: Sparse Autoencoder Pre-training}
An overcomplete SAE $(d \to D, D \gg d)$ learns sparse concepts via TopK activation:
\begin{align}
h &= \text{ReLU}(W_e(x - b_d) + b_e) \quad (\text{pre-activation}) \\
z &= \text{TopK}_k(h) \quad (\text{sparse concepts, $k \ll D$}) \\
\hat{x} &= W_d z + b_d \quad (\text{reconstruction})
\end{align}
Loss: $\mathcal{L}_\text{SAE} = \alpha \|x - \hat{x}\|^2 + \lambda_1 \|h\|_1 + \lambda_3 \mathcal{L}_\text{signal}(W_d, V_k)$

where $\mathcal{L}_\text{signal}$ penalizes decoder columns living in the noise subspace:
\begin{equation}
\mathcal{L}_\text{signal} = \frac{1}{D}\sum_{j=1}^D \left(1 - \|V_k^\top d_j\|^2\right)
\end{equation}
$d_j$ is the $j$-th decoder column (unit-norm enforced), and $V_k^\top d_j$ projects it onto the signal subspace.

\textbf{Key constraint:} Decoder columns $d_j$ are re-normalized to unit norm after each gradient step, ensuring $\|d_j\|_2 = 1$.

\paragraph{Fixed Normalization Buffer}
After SAE convergence, compute per-feature statistics \emph{from non-zero activations only}:
\begin{align}
\mu_j &= \mathbb{E}_{i: z_{ij} > 0}[z_{ij}] \\
\sigma_j &= \sqrt{\text{Var}_{i: z_{ij} > 0}[z_{ij}]}
\end{align}
Store as fixed buffers (not learnable). This normalization brings sparse activations (mean $\sim$180, max $\sim$1500) into the sigmoid's discriminative range $[-3, 3]$.

\paragraph{Stage 2B: Logic Layer Training (SAE Frozen)}
With SAE parameters frozen, train:

\textbf{1. Sigmoid Bottleneck:} Binarize concepts for logic operations
\begin{equation}
c_j = \sigma(\alpha_j(z_j' - \beta_j)), \quad z' = \frac{z - \mu_z}{\sigma_z}
\end{equation}
where $\alpha_j, \beta_j$ are learnable per-feature. A bimodality loss $\mathcal{L}_\text{bim} = \sum_j c_j(1-c_j)$ pushes outputs toward $\{0,1\}$.

\textbf{2. Product T-Norm Logic Layer:} For each action $a$, learn $n$ clauses in DNF form. A clause $C_i$ computes:
\begin{align}
\ell_{ij} &= p_{ij} \cdot c_j + n_{ij} \cdot (1 - c_j) + (1 - p_{ij} - n_{ij}) \quad \text{(soft literal)} \\
C_i &= \sigma\left(\sum_{j=1}^D \log(\ell_{ij} + \epsilon) + b_i\right) \quad \text{(clause activation)}
\end{align}
where $(p_{ij}, n_{ij}) = \text{softmax}(w_{ij}^{(+)}, w_{ij}^{(-)}; \text{dim}=0)$ encodes a 3-way choice:
\begin{itemize}
\item $p_{ij} \approx 1$: feature $j$ must be present ($f_j$)
\item $n_{ij} \approx 1$: feature $j$ must be absent ($\neg f_j$)
\item $p_{ij}, n_{ij} \approx 0$: feature $j$ ignored
\end{itemize}

The product t-norm (log-sum-exp) implements soft AND: $\prod_j \ell_{ij} \approx \exp(\sum_j \log \ell_{ij})$.

Action logits: $\text{score}_a = \sum_{i=1}^n C_{a,i}$ (OR of clauses).

\textbf{Training loss (Stage 2B only):}
\begin{align}
\mathcal{L}_\text{logic} = &\beta \cdot \text{CrossEntropy}(a, \hat{a}) \\
&+ \lambda_\text{bim}(t) \cdot \sum_j c_j(1-c_j) \\
&+ \lambda_0 \cdot \sum_{i,j} (p_{ij} + n_{ij})
\end{align}
where $\lambda_\text{bim}(t)$ is annealed from 0 to 0.3 over 80 epochs, and $\lambda_0$ encourages sparse rules.

\subsection{Training Protocol}

\begin{algorithm}[h]
\caption{Two-Stage ARC Training}
\begin{algorithmic}[1]
\State \textbf{Input:} Feature rollout $\{(x_i, a_i)\}_{i=1}^N$, hyperparameters
\State \textbf{Stage 1:} Preprocess
\State $\quad$ Run SVD on $X$ to get $V_k$, explained variance
\State $\quad$ Run ICA to get $k$ independent directions
\State $\quad$ Compute $\mu_x, \sigma_x$; normalize: $x_i \leftarrow (x_i - \mu_x)/\sigma_x$
\State \textbf{Stage 2A:} SAE Pre-training (50 epochs)
\State $\quad$ Initialize SAE decoder with ICA directions + random signal-subspace vectors
\State $\quad$ \textbf{for} $t = 1$ to $T_\text{SAE}$:
\State $\quad\quad$ Sample batch $\{x_i\}$
\State $\quad\quad$ Compute $z = \text{TopK}_k(\text{ReLU}(W_e(x - b_d) + b_e))$
\State $\quad\quad$ Compute $\hat{x} = W_d z + b_d$
\State $\quad\quad$ $\mathcal{L} = \|x - \hat{x}\|^2 + \lambda_1 \|h\|_1 + \lambda_3 \mathcal{L}_\text{signal}$
\State $\quad\quad$ Update $W_e, W_d, b_e, b_d$ via Adam
\State $\quad\quad$ Re-normalize decoder columns: $d_j \leftarrow d_j / \|d_j\|_2$
\State $\quad$ \textbf{Freeze} SAE parameters
\State $\quad$ Compute fixed $\mu_z, \sigma_z$ from non-zero $z$ activations
\State \textbf{Stage 2B:} Logic Training (350 epochs, SAE frozen)
\State $\quad$ \textbf{for} $t = 1$ to $T_\text{logic}$:
\State $\quad\quad$ Sample batch $\{x_i, a_i\}$
\State $\quad\quad$ $z = \text{SAE.encode}(x_i)$ \textcolor{gray}{\# frozen}
\State $\quad\quad$ $z' = (z - \mu_z)/\sigma_z$ \textcolor{gray}{\# fixed normalization}
\State $\quad\quad$ $c = \sigma(\alpha(z' - \beta))$ \textcolor{gray}{\# bottleneck}
\State $\quad\quad$ Compute action logits via product t-norm logic
\State $\quad\quad$ $\mathcal{L} = \text{CE}(a_i, \hat{a}_i) + \lambda_\text{bim}(t) \mathcal{L}_\text{bim} + \lambda_0 \mathcal{L}_\text{L0}$
\State $\quad\quad$ Update $\alpha, \beta, \{p_{ij}, n_{ij}\}, \{b_i\}$ via Adam
\State \textbf{Extract rules:} For each clause $C_i$, output $\{j: p_{ij} > 0.3\}$ as positive literals, $\{j: n_{ij} > 0.3\}$ as negative literals
\end{algorithmic}
\end{algorithm}

\textbf{Why two stages?} Joint training causes \emph{gradient wars}: SAE drifts to minimize reconstruction, continuously reshaping the feature space, which destroys the logic layer's learned mappings. Accuracy peaks early ($\sim$93\% at epoch 30) then collapses to 70\%. Freezing the SAE gives the logic layer a stationary target distribution.

\subsection{Rule Extraction}

After training, extract interpretable DNF rules by thresholding selector probabilities ($p_{ij}, n_{ij}$):
\begin{itemize}
\item If $p_{ij} > \tau$ (e.g., $\tau=0.3$): include positive literal $f_j$
\item If $n_{ij} > \tau$: include negative literal $\neg f_j$
\item Otherwise: feature $j$ ignored in clause $i$
\end{itemize}

For action $a$, the rule is:
\begin{equation}
a \leftarrow \bigvee_{i=1}^{n_a} \left(\bigwedge_{j: p_{ij}>\tau} f_j \land \bigwedge_{j: n_{ij}>\tau} \neg f_j\right)
\end{equation}

\section{Results}

\subsection{Experimental Setup}

\textbf{Environment:} MiniGrid-DoorKey-5x5-v0 \cite{minigrid}, a grid-world task requiring the agent to pick up a key, unlock a door, and reach a goal.

\textbf{Baseline Policy:} PPO-trained policy with CNN feature extractor ($d=128$), achieving 98\% success rate.

\textbf{Hyperparameters:} SAE hidden dimension $D=300$, TopK $k=50$, 10 clauses per action, 50 SAE pre-training epochs, 350 logic training epochs.

\subsection{Performance}

Evaluation over 100 episodes on MiniGrid-DoorKey-5x5:

\begin{table}[h]
\centering
\caption{Success rate comparison}
\begin{tabular}{lcc}
\toprule
Method & Success Rate & Interpretable \\
\midrule
PPO (original) & 98.0\% & No \\
ARC (ours) & 90--100\% & Yes \\
Linear Probe (upper bound) & 93.4\% & No \\
\bottomrule
\end{tabular}
\end{table}

The linear probe test (logistic regression on bottleneck features) achieves 93.4\% validation accuracy, indicating that the bottleneck features retain sufficient information. The logic layer's matching performance confirms that DNF rules are sufficiently expressive.

\subsection{Learned Rules}

Example extracted rules (bias terms indicate clause confidence):

\begin{verbatim}
TurnRight <-
  (~f_25 & ~f_26 & f_81)
  [bias=4.10, x10]

Forward <-
  (~f_25 & f_26 & ~f_81 & ~f_91 & ~f_280)
  [bias=4.91, x10]

Pickup <-
  (f_22 & f_25 & f_26 & ~f_81 & ~f_195 & f_255)
  [bias=5.55, x10]

Toggle <-
  (~f_22 & f_25 & ~f_26 & f_91 & f_195
   & ~f_255 & f_280)
  [bias=4.88, x10]
\end{verbatim}

Notation: $f_i$ is feature $i$, $\lnot f_i$ (or \texttt{\textasciitilde f\_i}) is its negation. ``x10'' indicates 10 redundant clauses were merged (learned same literals but different bias). Bias $b_i$ controls clause threshold.

\textbf{Rule Characteristics:}
\begin{itemize}
\item \textbf{Compact:} 4--7 literals per clause (vs. decision trees with 15--20 depth)
\item \textbf{Multi-variate:} Express conjunctions natively, not via cascaded binary splits
\item \textbf{Verifiable:} Can manually check each clause on test observations
\end{itemize}

\subsection{Concept Interpretability}

By visualizing maximally activating observations for each feature $f_i$, domain experts can assign semantic labels post-hoc. For instance:
\begin{itemize}
\item $f_{25}$: ``door in view''
\item $f_{81}$: ``key visible''
\item $f_{255}$: ``adjacent to object''
\end{itemize}

This yields human-readable rules like:
\begin{verbatim}
Pickup <-
  (key_visible & door_in_view & adjacent_to_object
   & ~obstacle_ahead & ...)
\end{verbatim}

\section{Discussion}

\subsection{Why ARC Works}

\paragraph{Separation of Concerns} The two-stage protocol disentangles \emph{what to represent} (SAE's job) from \emph{how to decide} (logic layer's job). This architectural constraint eliminates gradient conflicts that plagued end-to-end training.

\paragraph{Overcomplete Sparse Codes} An overcomplete dictionary ($D \gg d$) allows the SAE to discover more concepts than input dimensions, while TopK sparsity ensures only $k \ll D$ concepts activate per sample. This strikes a balance between expressiveness and interpretability.

\paragraph{Product T-Norm as Soft AND} The log-sum-exp approximation $\prod_j \ell_{ij} \approx \exp(\sum_j \log \ell_{ij})$ makes conjunctions differentiable. The 3-way softmax $(p, n, \text{absent})$ enables gradient-based feature selection within each clause.

\subsection{Comparison to Prior Work}

\begin{table}[h]
\centering
\caption{Interpretability method comparison}
\small
\begin{tabular}{lccc}
\toprule
Method & Manual & Accuracy & Rule \\
 & Concepts? & Drop? & Type \\
\midrule
CBM \cite{koh2020concept} & Yes & Moderate & Linear \\
Decision Tree & No & High & Hierarchical \\
Soft DT \cite{frosst2017distilling} & No & Moderate & Hierarchical \\
ARC (ours) & No & Low & DNF \\
\bottomrule
\end{tabular}
\end{table}

ARC uniquely combines automatic concept discovery with compact logical rules, avoiding both manual annotation and accuracy degradation.

\subsection{Limitations}

\paragraph{Environment Dependence} Current evaluation focuses on grid-world navigation. Scaling to high-dimensional visual domains (Atari, robotics) requires handling larger feature extractors ($d > 1000$) and more diverse concepts.

\paragraph{Concept Interpretability} While SAE features are sparse and visualizable, their semantic meaning requires post-hoc inspection. Incorporating auxiliary losses (e.g., matching human-provided concept labels when available) could improve alignment.

\paragraph{Rule Complexity} Although DNF rules are more compact than decision trees, clauses with 7+ literals may still challenge human comprehension. Future work could explore hierarchical rule organization or clause summarization.

\paragraph{Computational Cost} The two-stage protocol requires training the SAE to convergence before logic training begins, roughly doubling training time compared to end-to-end approaches. However, this is a one-time cost for improved stability.

\section{Conclusion}

We introduced ARC, a framework for extracting interpretable boolean rules from deep RL policies. By separating sparse concept learning from logical rule induction, ARC achieves 90--100\% success rate on MiniGrid while producing concise DNF rules. Our key insight—freezing the SAE before training the logic layer—eliminates gradient conflicts and stabilizes training.

Future directions include: (1) scaling to complex visual domains, (2) integrating weak supervision for concept grounding, (3) exploring hierarchical rule structures for deeper policies, and (4) applying ARC to safety-critical domains where verifiable interpretability is paramount.

\section*{Acknowledgments}
The authors thank the open-source contributors to Stable-Baselines3, MiniGrid, and PyTorch for enabling this research.

\begin{thebibliography}{9}

\bibitem{silver2016alphago}
D. Silver et al., ``Mastering the game of Go with deep neural networks and tree search,'' \emph{Nature}, vol. 529, pp. 484--489, 2016.

\bibitem{mnih2015human}
V. Mnih et al., ``Human-level control through deep reinforcement learning,'' \emph{Nature}, vol. 518, pp. 529--533, 2015.

\bibitem{koh2020concept}
P. W. Koh et al., ``Concept bottleneck models,'' in \emph{ICML}, 2020.

\bibitem{breiman1984classification}
L. Breiman et al., \emph{Classification and Regression Trees}, Wadsworth, 1984.

\bibitem{frosst2017distilling}
N. Frosst and G. Hinton, ``Distilling a neural network into a soft decision tree,'' in \emph{CEUR Workshop Proceedings}, 2017.

\bibitem{minigrid}
M. Chevalier-Boisvert et al., ``Minimalistic gridworld environment for OpenAI Gym,'' \url{https://github.com/Farama-Foundation/Minigrid}, 2018.

\end{thebibliography}

\end{document}
